\documentclass[11pt]{article}
\usepackage[margin=1in]{geometry}
\usepackage{amsmath,amssymb,amsthm}
\title{Disproof of Conjecture 00000001068\\
\large The maximal sum-free subsets of $\mathbb{F}_p$ are \emph{not} exactly the intervals $((p+1)/3,\,2(p-1)/3)$ up to dilation}
\author{TLMC submission}
\date{October 2, 2026}

\newcommand{\F}{\mathbb{F}}
\newtheorem*{claim}{Conjecture 00000001068 (as stated)}
\newtheorem*{thm}{Theorem}

\begin{document}
\maketitle

\section{The conjecture}

\begin{claim}
Definition: a sum-free set is a set $A$ with $A + A$ disjoint from $A$. The maximal sum-free subsets of $\F_p$ are exactly the intervals $((p+1)/3,\, 2(p-1)/3)$, uniquely up to dilation (the complete Diananda--Yap characterization over finite fields).
\end{claim}

\section{Counterexample: $p = 11$}

Let $A = \{4,5,6,7\} \subseteq \F_{11}$.

\begin{thm}
$A$ is an inclusion-maximal sum-free subset of $\F_{11}$, but $A$ is not a dilation of $((p+1)/3,\,2(p-1)/3) = (4,\,20/3)$. Hence the conjecture is false.
\end{thm}

\begin{proof}
\emph{Sum-free.} Computing all pairwise sums in $\F_{11}$:
\[
A + A \;=\; \{8,9,10,0,\; 9,10,0,1,\; 10,0,1,2,\; 0,1,2,3\} \;=\; \{0,1,2,3,8,9,10\} \;=\; \F_{11}\setminus A,
\]
which is disjoint from $A$. (So $A$ is even a \emph{maximum} sum-free set: $|A| = 4$, and no sum-free set in $\F_{11}$ contains a proper superset of $A$.)

\emph{Inclusion-maximal.} Every $x \in \F_{11}\setminus A$ lies in $A+A$; hence $x \in (A\cup\{x\}) + (A\cup\{x\})$ and $A\cup\{x\}$ is not sum-free. So $A$ is an inclusion-maximal sum-free subset of $\F_{11}$.

\emph{The conjectured family.} With rational (open) interval endpoints,
\[
\bigl((p+1)/3,\; 2(p-1)/3\bigr) \;=\; \bigl(4,\; 20/3\bigr) \cap \mathbb{Z} \;=\; \{5,6\} \;\subseteq\; \F_{11},
\]
which has only $|I| = 2$ elements. Every dilation $cI$ with $c \in \F_{11}^{\times}$ again has exactly $2$ elements, and $0\cdot I = \{0\}$. Exhaustively, over all $c \in \{0,1,\dots,10\}$ one has $c \cdot \{5,6\} \neq \{4,5,6,7\}$ (the distinct images are $\{0\},\{5,6\},\{10,1\},\{4,7\},\{9,2\},\{3,8\}$).

Since $|A| = 4 \neq 2$, $A$ is an inclusion-maximal sum-free subset of $\F_{11}$ that is not of the form $c\bigl((p+1)/3,\,2(p-1)/3\bigr)$. The claimed characterization fails.
\end{proof}

\section{Remarks}

The failure is the open-interval formula itself. The classical Diananda--Yap extremal family in $\mathbb{Z}_p$ is the \emph{closed} integer interval
$\{\lfloor p/3\rfloor + 1, \dots, \lfloor 2p/3 \rfloor\}$,
which at $p=11$ is $\{4,5,6,7\} = A$ itself (size $\lfloor 2p/3\rfloor - \lfloor p/3\rfloor = 4$); the conjecture's open interval $(4, 20/3)$ drops the two endpoints and keeps only $2$ points, so the family it defines is not maximal.

\section{Machine verification}

\begin{itemize}
\item \texttt{reproduce.py} recomputes $A+A$, the interval, all dilations, sum-freeness, and inclusion-maximality from scratch; all checks pass.
\item \texttt{lean4/Main.lean} formalizes the counterexample in Lean 4 (v4.33.1, core only): \texttt{A\_sum\_free}, \texttt{A\_maximal}, \texttt{A\_len}, and \texttt{A\_not\_dilation} combine into \texttt{counterexample\_1068}. All proofs are \texttt{rfl} over explicit $\F_{11}$ computations. The audit in \texttt{Check.lean} reports \textbf{7/7 theorems depend on no axioms} (no \texttt{sorry}, no \texttt{propext}, no \texttt{Classical.choice}, no \texttt{Quot.sound}).
\end{itemize}

\end{document}
